% =====================================================================
%  rw_related_work.tex  --  Chapter 3: Related Work
%  WHERE:  replaces the current content of Chapter 3 (the old three-line
%          paragraph about sentiment analysis / Word2Vec / GloVe / FastText,
%          which belongs to the previous template, not to this thesis).
%          Use:   % =====================================================================
%  rw_related_work.tex  --  Chapter 3: Related Work
%  WHERE:  replaces the current content of Chapter 3 (the old three-line
%          paragraph about sentiment analysis / Word2Vec / GloVe / FastText,
%          which belongs to the previous template, not to this thesis).
%          Use:   % =====================================================================
%  rw_related_work.tex  --  Chapter 3: Related Work
%  WHERE:  replaces the current content of Chapter 3 (the old three-line
%          paragraph about sentiment analysis / Word2Vec / GloVe / FastText,
%          which belongs to the previous template, not to this thesis).
%          Use:   % =====================================================================
%  rw_related_work.tex  --  Chapter 3: Related Work
%  WHERE:  replaces the current content of Chapter 3 (the old three-line
%          paragraph about sentiment analysis / Word2Vec / GloVe / FastText,
%          which belongs to the previous template, not to this thesis).
%          Use:   \input{related_work/rw_related_work}
%  NEEDS:  bg_preamble.tex, rw_preamble.tex, bg_references.bib,
%          rw_references.bib
% =====================================================================

\chapter{Related Work}
\label{chap:related-work}

This chapter reviews the research that this thesis builds on and
distinguishes it from existing approaches. The thesis lies at the
intersection of three research areas (Figure~\ref{fig:rw-positioning}):
(i)~large language models for medicine and patient education,
(ii)~knowledge graphs and neuro-symbolic reasoning for LLMs, and
(iii)~automatic prompt optimisation and the automatic generation of
educational videos. Section~\ref{sec:rw-medical-llm} discusses medical LLMs
and their limitations for patient education.
Section~\ref{sec:rw-kg-llm} reviews approaches that ground medical LLMs in
knowledge graphs, and Section~\ref{sec:rw-nesy} reviews methods that combine
symbolic rules with LLM reasoning. Section~\ref{sec:rw-prompt} covers
automatic prompt optimisation, and Section~\ref{sec:rw-video} covers the
generation of educational videos. Section~\ref{sec:rw-summary} compares all
approaches and summarises the research gap addressed by this thesis.

\begin{figure}[htbp]
  \centering
  \begin{tikzpicture}[font=\footnotesize]
    \def\r{3.0}
    \begin{scope}[blend group=multiply]
      \fill[bgOrangeLight] (0,0) circle (\r);
      \fill[bgBlueLight]   (3.4,0) circle (\r);
      \fill[bgPurpleLight] (1.7,-3.0) circle (\r);
    \end{scope}
    \draw[bgOrange, thick] (0,0) circle (\r);
    \draw[bgBlue, thick]   (3.4,0) circle (\r);
    \draw[bgPurple, thick] (1.7,-3.0) circle (\r);
    % area titles (outside the circles)
    \node[align=center, text=bgOrange, font=\footnotesize\bfseries] at (-1.4,3.45) {LLMs for medicine\\\& patient education};
    \node[align=center, text=bgBlue, font=\footnotesize\bfseries]   at (4.8,3.45) {Knowledge graphs \&\\neuro-symbolic AI};
    \node[align=center, text=bgPurple, font=\footnotesize\bfseries] at (1.7,-6.5) {Prompt optimisation \&\\educational video generation};
    % works that belong to one area only
    \node[align=center, font=\scriptsize] at (-1.25,0.1) {Med-PaLM~\cite{singhal2023large}\\Med-HALT~\cite{pal2023medhalt}\\King et al.~\cite{king2025improving}};
    \node[align=center, font=\scriptsize] at (4.65,0.1) {AMIE~3~\cite{lajus2020fast}\\RoG~\cite{luo2024reasoning}\\ChatRule~\cite{luo2025chatrule}\\BRINK~\cite{zhou2026breaks}};
    \node[align=center, font=\scriptsize] at (1.7,-4.35) {GAAPO~\cite{secheresse2025gaapo}\\OPRO~\cite{yang2024large}\\Code2Video~\cite{chen2026code2video}};
    % overlaps
    \node[align=center, font=\tiny] at (1.7,0.85) {KG-RAG~\cite{soman2024biomedical}\\MedGraphRAG~\cite{wu2025medical}\\DR.KNOWS~\cite{gao2025leveraging}};
    \node[align=center, font=\tiny] at (-0.1,-2.05) {Medprompt\\\cite{nori2023can}};
    % centre
    \node[draw=black, thick, fill=white, rounded corners=3pt, align=center,
          font=\scriptsize\bfseries, inner sep=3pt] at (1.7,-0.95) {This\\thesis};
  \end{tikzpicture}
  \caption[Positioning of this thesis within related research areas]{Positioning
  of this thesis. Existing work covers one or two of the three research areas;
  this thesis combines guideline-based knowledge graphs, rule-guided
  reasoning, prompt optimisation and patient-oriented storyboards for
  educational videos.}
  \label{fig:rw-positioning}
\end{figure}

% ---------------------------------------------------------------------
\section{Large Language Models in Medicine and Patient Education}
\label{sec:rw-medical-llm}
% ---------------------------------------------------------------------

\paragraph{Medical question answering.}
Singhal et al.~\cite{singhal2023large} introduced MultiMedQA, a benchmark that
combines professional medical exams, research questions and consumer health
questions, and adapted the general-purpose model Flan-PaLM to the medical
domain through instruction prompt tuning. In a human evaluation, 92.6\,\% of
the long-form answers of the resulting model, Med-PaLM, were judged to be in
line with scientific consensus, compared with 92.9\,\% of clinician answers;
nevertheless, the model remained inferior to clinicians on several other
evaluation axes. Its successor
Med-PaLM~2~\cite{singhal2025toward} further narrowed this gap. Nori et
al.~\cite{nori2023can} showed that a general model (GPT-4) can reach similar
results without any medical fine-tuning when it is combined with a carefully
designed prompting strategy (Section~\ref{sec:rw-prompt}). Surveys of the
field~\cite{thirunavukarasu2023large} conclude that LLMs are promising for
clinical documentation, education and patient communication, but that their
reliability, transparency and accountability are not yet sufficient for
unsupervised use.

\paragraph{Patient communication and patient education.}
LLMs have also been evaluated directly on patient-facing tasks. Ayers et
al.~\cite{ayers2023comparing} compared answers by physicians and by ChatGPT to
195 patient questions from a public online forum; in 78.6\,\% of 585
evaluations, healthcare professionals preferred the chatbot answer, which was
rated higher for both quality and empathy. In cardiology, King et
al.~\cite{king2025improving} used GPT-4 to simplify 143 heart-failure patient
education materials from leading US cardiology institutions. The median
reading level decreased from grade~10.3 to grade~7.3, and a cardiologist rated
all simplified texts as equally accurate as the originals. These studies show
that LLMs can make medical information more accessible. However, they rely
on general-purpose models without explicit medical knowledge, they evaluate
plain text only, and correctness is checked manually after generation.

\paragraph{Hallucinations in medical LLMs.}
The main obstacle to the safe use of LLMs for patients is hallucination
(Section~\ref{sec:bg-hallucination}). Pal et al.~\cite{pal2023medhalt}
proposed Med-HALT, a benchmark that tests both reasoning hallucinations (e.g.\
choosing an answer to a question with a false premise) and memory-based
hallucinations (e.g.\ inventing references), and found that all evaluated
LLMs were affected. Kim et al.~\cite{kim2025medicalhallucinationsfoundationmodels}
analysed eleven foundation models on seven medical tasks and, together with a
survey of clinicians, concluded that medical hallucinations are frequent in
practice. Interestingly, they also observed that chain-of-thought prompting
reduced hallucinations in most comparisons, which supports the idea of making
the reasoning of a model explicit.

\begin{rwgap}
Medical LLMs can produce fluent and empathetic answers, but they are not
explicitly grounded in clinical guidelines, their reasoning is not
verifiable, and existing patient-education studies focus on text.
Approaches are needed that connect patient-oriented generation with
explicit, guideline-based medical knowledge.
\end{rwgap}

% ---------------------------------------------------------------------
\section{Knowledge-Graph-Enhanced LLMs in the Biomedical Domain}
\label{sec:rw-kg-llm}
% ---------------------------------------------------------------------

A common way to reduce hallucinations is to supply LLMs with external
knowledge at inference time (Section~\ref{sec:bg-rag}). Xiong et
al.~\cite{xiong2024benchmarking} systematically benchmarked
retrieval-augmented generation for medicine on MIRAGE, a collection of 7,663
questions from five medical QA datasets. Their MedRAG toolkit improved the
accuracy of six LLMs by up to 18\,\% compared with chain-of-thought prompting.
These systems retrieve unstructured text passages, for example from PubMed or
textbooks.

Several works retrieve structured knowledge from \emph{knowledge graphs}
instead. Soman et al.~\cite{soman2024biomedical} proposed KG-RAG, which
extracts disease-specific context from SPOKE, a large biomedical knowledge
graph, and adds it to the prompt. By pruning the retrieved context with
embeddings, they reduced the number of prompt tokens by more than 50\,\%
while improving the answers of models such as Llama-2 and GPT-4. Gao et
al.~\cite{gao2025leveraging} presented DR.KNOWS, which selects paths from a
UMLS-based knowledge graph that are relevant to a patient's electronic health
record and uses them to support diagnosis prediction; clinicians judged the
resulting diagnostic rationales to be clinically sound. Wu et
al.~\cite{wu2025medical} proposed MedGraphRAG, which links private medical
documents to trusted medical sources and vocabularies in a three-level graph,
and retrieves from this graph so that answers are supported by evidence.
Outside medicine, Think-on-Graph~\cite{sun2024thinkongraph} lets an LLM
explore a knowledge graph step by step with beam search, and Microsoft's
GraphRAG~\cite{edge2024local} builds a graph of entities and communities from
text to answer global questions about a document collection.

All of these approaches use the knowledge graph only at \emph{inference
time} and assume that the facts needed to answer a question are present in
the graph. Zhou et al.~\cite{zhou2026breaks} questioned this assumption with
the BRINK benchmark (Section~\ref{sec:bg-rag}): when the triple that directly
supports an answer is removed, current KG-RAG methods show limited reasoning
ability, often fall back on knowledge memorised during pre-training, and
generalise differently depending on their design. This finding is highly
relevant for medical knowledge graphs, which are incomplete by nature
(Section~\ref{sec:bg-owa}).

\begin{bgexample}[Retrieval versus learnt reasoning under missing knowledge]
Assume the guideline knowledge graph contains the facts
\emph{(Drug~X, belongsToClass, MRA)} and
\emph{(MRA, recommendedFor, HFrEF)}, but \emph{not} the direct fact
\emph{(Drug~X, recommendedFor, HFrEF)}. A patient asks:
\emph{``Is Drug~X recommended for my type of heart failure?''}
\begin{itemize}
  \item A \textbf{retrieval-only} system searches for the direct triple, does
        not find it, and either answers ``no information'' or falls back on
        what the LLM memorised during pre-training.
  \item A model that was \textbf{fine-tuned on rule-based reasoning chains},
        e.g.\ derived from the rule
        $\mathit{belongsToClass}(x,c) \wedge \mathit{recommendedFor}(c,d)
        \Rightarrow \mathit{recommendedFor}(x,d)$,
        has learnt to combine the two remaining facts and can answer:
        \emph{``Drug~X is an MRA, and MRAs are recommended for HFrEF;
        therefore \ldots''}, while showing the facts it used.
\end{itemize}
This illustrates why this thesis combines knowledge-graph context at
inference time with reasoning patterns learnt from symbolic rules.
\end{bgexample}

\begin{rwgap}
Biomedical KG-enhanced LLMs ground answers in retrieved facts but do not
teach the model how to \emph{reason} over the graph when facts are missing.
Moreover, they target clinicians or benchmark questions, not patient
education based on clinical practice guidelines and patient questionnaires.
\end{rwgap}

% ---------------------------------------------------------------------
\section{Symbolic Rules and Neuro-Symbolic Reasoning with LLMs}
\label{sec:rw-nesy}
% ---------------------------------------------------------------------

\paragraph{Rule learning over knowledge graphs.}
Learning logical rules from knowledge graphs has a long tradition. The AMIE
family~\cite{galarraga2013amie,galarraga2015fast,lajus2020fast} mines closed
Horn rules under the Partial Completeness Assumption and is the rule miner
used in this thesis (Section~\ref{sec:bg-rule-mining}). AnyBURL~\cite{meilicke2019anytime}
learns rules bottom-up by sampling paths from the graph and generalising them,
which makes it efficient for knowledge-graph completion. Neural approaches
such as Neural~LP~\cite{yang2017differentiable} learn rules in a
differentiable way, trading some interpretability for end-to-end training.
In all of these works, rules are used to \emph{complete} the knowledge graph,
not to train a language model.

\paragraph{LLMs and rules.}
More recent work connects rules with LLMs. ChatRule~\cite{luo2025chatrule}
uses an LLM to \emph{generate} candidate rules from the semantic and
structural information of a knowledge graph and ranks them with
measures computed on the existing facts. Reasoning on Graphs
(RoG)~\cite{luo2024reasoning} fine-tunes an LLM to first generate relation
paths that are grounded in the knowledge graph (``planning'') and then to
reason over the retrieved paths, producing faithful and interpretable
answers. Closest to this thesis is Chain-of-Knowledge~\cite{zhang2024chain},
which mines rules from knowledge graphs, turns them into a knowledge-reasoning
dataset (KnowReason), and fine-tunes LLMs on it; to prevent the model from
simply memorising the rules, the authors add a trial-and-error mechanism.

\paragraph{Reasoning chains as training data.}
Training models on explicit reasoning chains has been shown to be
effective in general. Hsieh et al.~\cite{hsieh2023distilling} demonstrated
that smaller models fine-tuned on rationales extracted from a large LLM can
outperform much larger models while using less training data. Rule
groundings from a knowledge graph provide such rationales without a teacher
LLM, and their correctness can be traced back to explicit facts
(Section~\ref{sec:bg-rules-to-cot}).

\begin{rwgap}
Neuro-symbolic approaches that combine rules and LLMs have so far been
evaluated on open-domain benchmark knowledge graphs (e.g.\ Freebase-based
question-answering datasets).
It is unclear whether rule-guided chain-of-thought fine-tuning improves LLM
performance in the medical domain, where knowledge comes from heterogeneous
sources (guidelines, terminologies and patient questionnaires) and where
rules must be filtered by confidence before they are used.
\end{rwgap}

% ---------------------------------------------------------------------
\section{Automatic Prompt Optimisation}
\label{sec:rw-prompt}
% ---------------------------------------------------------------------

Since the quality of LLM outputs depends strongly on the prompt, several
methods search for better prompts automatically. Automatic Prompt Engineer
(APE)~\cite{zhou2023large} lets an LLM propose candidate instructions from
input--output examples and selects the best one by its score. ProTeGi
(APO)~\cite{pryzant2023automatic} improves prompts with natural-language
``gradients'', and OPRO~\cite{yang2024large} uses an LLM as an optimiser that
sees the history of previous prompts and their scores. Evolutionary
approaches maintain a population of prompts: EvoPrompt~\cite{guo2024connecting}
uses an LLM to perform mutation and crossover and improved over
human-written prompts on 31 datasets, and Promptbreeder~\cite{fernando2024promptbreeder}
additionally evolves the mutation prompts themselves. GAAPO~\cite{secheresse2025gaapo}
combines mutation, crossover, few-shot augmentation, APO and OPRO in a single
genetic algorithm and records which generator produced each improvement
(Section~\ref{sec:bg-apo}).

In medicine, Medprompt~\cite{nori2023can} combined dynamic few-shot
selection, self-generated chain-of-thought explanations and choice-shuffle
ensembling, which allowed GPT-4 to reach 90.2\,\% accuracy on the MedQA
benchmark without any domain-specific fine-tuning. Medprompt is a fixed
composition of prompting techniques, not a search over prompts, and like most
prompt-optimisation methods it was evaluated on multiple-choice questions.

\begin{rwgap}
Prompt-optimisation methods are mostly evaluated on classification and
multiple-choice tasks with very large models. They have not been studied for
generating \emph{structured} patient-education content, such as video
storyboards, with a small, domain-fine-tuned model whose medical content must
not be changed by the optimised prompt.
\end{rwgap}

% ---------------------------------------------------------------------
\section{Automatic Generation of Educational Videos}
\label{sec:rw-video}
% ---------------------------------------------------------------------

Video is an effective format for patient education: a systematic review
found strong evidence that video-based interventions improve the knowledge
of patients with chronic illnesses~\cite{deshpande2023video}. Producing such
videos, however, requires medical experts, script writers and designers.
Recent work therefore investigates how LLMs can automate parts of this
process. TheoremExplainAgent~\cite{ku2025theoremexplainagent} uses LLM agents
to plan and write Manim animation code for explanation videos of
mathematical and scientific theorems; the best configuration produced
videos for 93.8\,\% of 240 theorems, although most videos still contained
minor visual layout problems. Code2Video~\cite{chen2026code2video} also
generates educational videos as executable code, using a planner, a coder and
a critic agent, and reports a 40\,\% improvement over direct code generation.
In healthcare, Haider et al.~\cite{haider2026ai} compared AI-generated avatar
videos with text handouts for postoperative education in a pilot randomised
trial. Knowledge scores did not differ significantly, but participants rated
the videos as clearer and more engaging.

These systems focus on the \emph{production} of a video from a given topic
and rely on the general knowledge of the LLM for its content. They do not
verify the medical statements against a trusted knowledge source, and they do
not adapt the content to the needs of a specific patient population.

\begin{rwgap}
LLM-based video generation has been studied mainly for STEM topics, while
medical studies use generic AI tools without knowledge grounding. A method is
missing that produces patient-oriented video storyboards whose medical
content is derived from guideline-based knowledge and whose structure can be
validated before production.
\end{rwgap}

% ---------------------------------------------------------------------
\section{Summary and Research Gap}
\label{sec:rw-summary}
% ---------------------------------------------------------------------

Table~\ref{tab:rw-comparison} compares representative approaches along the
dimensions that are relevant for this thesis. Each line of research covers
only part of the problem: medical LLMs and patient-education studies provide
patient-oriented language but no explicit knowledge; biomedical KG-RAG
systems ground answers in knowledge graphs but only at inference time;
neuro-symbolic methods learn from rules but have not been applied to
medicine; prompt optimisation and video generation have not been combined
with knowledge-grounded medical content.

\begin{table}[tbp]
  \centering
  \small
  \setlength{\tabcolsep}{4pt}
  \begin{tabular}{@{}lccccccc@{}}
    \toprule
    \textbf{Approach} & \textbf{Med.} & \textbf{KG} & \textbf{Rules}
      & \textbf{CoT-FT} & \textbf{P-Opt} & \textbf{Patient} & \textbf{Video} \\
    \midrule
    Med-PaLM~\cite{singhal2023large}                 & \rwyes & \rwno  & \rwno  & \rwno  & \rwno  & \rwpart & \rwno \\
    King et al.~\cite{king2025improving}             & \rwyes & \rwno  & \rwno  & \rwno  & \rwno  & \rwyes  & \rwno \\
    MedRAG~\cite{xiong2024benchmarking}              & \rwyes & \rwno  & \rwno  & \rwno  & \rwno  & \rwno   & \rwno \\
    KG-RAG~\cite{soman2024biomedical}                & \rwyes & \rwyes & \rwno  & \rwno  & \rwno  & \rwno   & \rwno \\
    DR.KNOWS~\cite{gao2025leveraging}                & \rwyes & \rwyes & \rwno  & \rwno  & \rwno  & \rwno   & \rwno \\
    MedGraphRAG~\cite{wu2025medical}                 & \rwyes & \rwyes & \rwno  & \rwno  & \rwno  & \rwno   & \rwno \\
    RoG~\cite{luo2024reasoning}                      & \rwno  & \rwyes & \rwpart& \rwpart& \rwno  & \rwno   & \rwno \\
    ChatRule~\cite{luo2025chatrule}                  & \rwno  & \rwyes & \rwyes & \rwno  & \rwno  & \rwno   & \rwno \\
    Chain-of-Knowledge~\cite{zhang2024chain}         & \rwno  & \rwyes & \rwyes & \rwyes & \rwno  & \rwno   & \rwno \\
    Medprompt~\cite{nori2023can}                     & \rwyes & \rwno  & \rwno  & \rwno  & \rwpart& \rwno   & \rwno \\
    GAAPO~\cite{secheresse2025gaapo}                 & \rwno  & \rwno  & \rwno  & \rwno  & \rwyes & \rwno   & \rwno \\
    TheoremExplainAgent~\cite{ku2025theoremexplainagent} & \rwno & \rwno & \rwno & \rwno & \rwno & \rwno   & \rwyes \\
    Code2Video~\cite{chen2026code2video}             & \rwno  & \rwno  & \rwno  & \rwno  & \rwno  & \rwno   & \rwyes \\
    \midrule
    \textbf{This thesis}                             & \rwyes & \rwyes & \rwyes & \rwyes & \rwyes & \rwyes  & \rwyes \\
    \bottomrule
  \end{tabular}
  \caption[Comparison of related approaches]{Comparison of related
  approaches. \textbf{Med.}: medical domain; \textbf{KG}: uses a knowledge
  graph; \textbf{Rules}: uses symbolic rules; \textbf{CoT-FT}: fine-tunes the
  LLM on (rule-based) reasoning chains; \textbf{P-Opt}: automatic prompt
  optimisation; \textbf{Patient}: patient-oriented output; \textbf{Video}:
  educational video or storyboard generation. \rwyes~=~yes,
  \rwpart~=~partially, \rwno~=~no. For this thesis, \textbf{Video} refers to
  storyboard scripts as input for video production.}
  \label{tab:rw-comparison}
\end{table}

Based on this analysis, the research gap addressed in this thesis can be
summarised as follows.

\begin{rwgap}[Research Gap of this Thesis]
There is no approach that (i)~represents clinical guidelines and patient
questionnaires about cardiovascular disease as knowledge graphs,
(ii)~turns symbolic rules mined from these graphs into chain-of-thought
training data for a medical LLM, (iii)~evaluates whether this improves
reasoning when knowledge is incomplete, and (iv)~uses the resulting model,
together with an automatically optimised prompt, to generate
patient-oriented storyboards for educational videos. This thesis addresses
this gap with the semantic-driven approach presented in
Chapter~4.
\end{rwgap}

%  NEEDS:  bg_preamble.tex, rw_preamble.tex, bg_references.bib,
%          rw_references.bib
% =====================================================================

\chapter{Related Work}
\label{chap:related-work}

This chapter reviews the research that this thesis builds on and
distinguishes it from existing approaches. The thesis lies at the
intersection of three research areas (Figure~\ref{fig:rw-positioning}):
(i)~large language models for medicine and patient education,
(ii)~knowledge graphs and neuro-symbolic reasoning for LLMs, and
(iii)~automatic prompt optimisation and the automatic generation of
educational videos. Section~\ref{sec:rw-medical-llm} discusses medical LLMs
and their limitations for patient education.
Section~\ref{sec:rw-kg-llm} reviews approaches that ground medical LLMs in
knowledge graphs, and Section~\ref{sec:rw-nesy} reviews methods that combine
symbolic rules with LLM reasoning. Section~\ref{sec:rw-prompt} covers
automatic prompt optimisation, and Section~\ref{sec:rw-video} covers the
generation of educational videos. Section~\ref{sec:rw-summary} compares all
approaches and summarises the research gap addressed by this thesis.

\begin{figure}[htbp]
  \centering
  \begin{tikzpicture}[font=\footnotesize]
    \def\r{3.0}
    \begin{scope}[blend group=multiply]
      \fill[bgOrangeLight] (0,0) circle (\r);
      \fill[bgBlueLight]   (3.4,0) circle (\r);
      \fill[bgPurpleLight] (1.7,-3.0) circle (\r);
    \end{scope}
    \draw[bgOrange, thick] (0,0) circle (\r);
    \draw[bgBlue, thick]   (3.4,0) circle (\r);
    \draw[bgPurple, thick] (1.7,-3.0) circle (\r);
    % area titles (outside the circles)
    \node[align=center, text=bgOrange, font=\footnotesize\bfseries] at (-1.4,3.45) {LLMs for medicine\\\& patient education};
    \node[align=center, text=bgBlue, font=\footnotesize\bfseries]   at (4.8,3.45) {Knowledge graphs \&\\neuro-symbolic AI};
    \node[align=center, text=bgPurple, font=\footnotesize\bfseries] at (1.7,-6.5) {Prompt optimisation \&\\educational video generation};
    % works that belong to one area only
    \node[align=center, font=\scriptsize] at (-1.25,0.1) {Med-PaLM~\cite{singhal2023large}\\Med-HALT~\cite{pal2023medhalt}\\King et al.~\cite{king2025improving}};
    \node[align=center, font=\scriptsize] at (4.65,0.1) {AMIE~3~\cite{lajus2020fast}\\RoG~\cite{luo2024reasoning}\\ChatRule~\cite{luo2025chatrule}\\BRINK~\cite{zhou2026breaks}};
    \node[align=center, font=\scriptsize] at (1.7,-4.35) {GAAPO~\cite{secheresse2025gaapo}\\OPRO~\cite{yang2024large}\\Code2Video~\cite{chen2026code2video}};
    % overlaps
    \node[align=center, font=\tiny] at (1.7,0.85) {KG-RAG~\cite{soman2024biomedical}\\MedGraphRAG~\cite{wu2025medical}\\DR.KNOWS~\cite{gao2025leveraging}};
    \node[align=center, font=\tiny] at (-0.1,-2.05) {Medprompt\\\cite{nori2023can}};
    % centre
    \node[draw=black, thick, fill=white, rounded corners=3pt, align=center,
          font=\scriptsize\bfseries, inner sep=3pt] at (1.7,-0.95) {This\\thesis};
  \end{tikzpicture}
  \caption[Positioning of this thesis within related research areas]{Positioning
  of this thesis. Existing work covers one or two of the three research areas;
  this thesis combines guideline-based knowledge graphs, rule-guided
  reasoning, prompt optimisation and patient-oriented storyboards for
  educational videos.}
  \label{fig:rw-positioning}
\end{figure}

% ---------------------------------------------------------------------
\section{Large Language Models in Medicine and Patient Education}
\label{sec:rw-medical-llm}
% ---------------------------------------------------------------------

\paragraph{Medical question answering.}
Singhal et al.~\cite{singhal2023large} introduced MultiMedQA, a benchmark that
combines professional medical exams, research questions and consumer health
questions, and adapted the general-purpose model Flan-PaLM to the medical
domain through instruction prompt tuning. In a human evaluation, 92.6\,\% of
the long-form answers of the resulting model, Med-PaLM, were judged to be in
line with scientific consensus, compared with 92.9\,\% of clinician answers;
nevertheless, the model remained inferior to clinicians on several other
evaluation axes. Its successor
Med-PaLM~2~\cite{singhal2025toward} further narrowed this gap. Nori et
al.~\cite{nori2023can} showed that a general model (GPT-4) can reach similar
results without any medical fine-tuning when it is combined with a carefully
designed prompting strategy (Section~\ref{sec:rw-prompt}). Surveys of the
field~\cite{thirunavukarasu2023large} conclude that LLMs are promising for
clinical documentation, education and patient communication, but that their
reliability, transparency and accountability are not yet sufficient for
unsupervised use.

\paragraph{Patient communication and patient education.}
LLMs have also been evaluated directly on patient-facing tasks. Ayers et
al.~\cite{ayers2023comparing} compared answers by physicians and by ChatGPT to
195 patient questions from a public online forum; in 78.6\,\% of 585
evaluations, healthcare professionals preferred the chatbot answer, which was
rated higher for both quality and empathy. In cardiology, King et
al.~\cite{king2025improving} used GPT-4 to simplify 143 heart-failure patient
education materials from leading US cardiology institutions. The median
reading level decreased from grade~10.3 to grade~7.3, and a cardiologist rated
all simplified texts as equally accurate as the originals. These studies show
that LLMs can make medical information more accessible. However, they rely
on general-purpose models without explicit medical knowledge, they evaluate
plain text only, and correctness is checked manually after generation.

\paragraph{Hallucinations in medical LLMs.}
The main obstacle to the safe use of LLMs for patients is hallucination
(Section~\ref{sec:bg-hallucination}). Pal et al.~\cite{pal2023medhalt}
proposed Med-HALT, a benchmark that tests both reasoning hallucinations (e.g.\
choosing an answer to a question with a false premise) and memory-based
hallucinations (e.g.\ inventing references), and found that all evaluated
LLMs were affected. Kim et al.~\cite{kim2025medicalhallucinationsfoundationmodels}
analysed eleven foundation models on seven medical tasks and, together with a
survey of clinicians, concluded that medical hallucinations are frequent in
practice. Interestingly, they also observed that chain-of-thought prompting
reduced hallucinations in most comparisons, which supports the idea of making
the reasoning of a model explicit.

\begin{rwgap}
Medical LLMs can produce fluent and empathetic answers, but they are not
explicitly grounded in clinical guidelines, their reasoning is not
verifiable, and existing patient-education studies focus on text.
Approaches are needed that connect patient-oriented generation with
explicit, guideline-based medical knowledge.
\end{rwgap}

% ---------------------------------------------------------------------
\section{Knowledge-Graph-Enhanced LLMs in the Biomedical Domain}
\label{sec:rw-kg-llm}
% ---------------------------------------------------------------------

A common way to reduce hallucinations is to supply LLMs with external
knowledge at inference time (Section~\ref{sec:bg-rag}). Xiong et
al.~\cite{xiong2024benchmarking} systematically benchmarked
retrieval-augmented generation for medicine on MIRAGE, a collection of 7,663
questions from five medical QA datasets. Their MedRAG toolkit improved the
accuracy of six LLMs by up to 18\,\% compared with chain-of-thought prompting.
These systems retrieve unstructured text passages, for example from PubMed or
textbooks.

Several works retrieve structured knowledge from \emph{knowledge graphs}
instead. Soman et al.~\cite{soman2024biomedical} proposed KG-RAG, which
extracts disease-specific context from SPOKE, a large biomedical knowledge
graph, and adds it to the prompt. By pruning the retrieved context with
embeddings, they reduced the number of prompt tokens by more than 50\,\%
while improving the answers of models such as Llama-2 and GPT-4. Gao et
al.~\cite{gao2025leveraging} presented DR.KNOWS, which selects paths from a
UMLS-based knowledge graph that are relevant to a patient's electronic health
record and uses them to support diagnosis prediction; clinicians judged the
resulting diagnostic rationales to be clinically sound. Wu et
al.~\cite{wu2025medical} proposed MedGraphRAG, which links private medical
documents to trusted medical sources and vocabularies in a three-level graph,
and retrieves from this graph so that answers are supported by evidence.
Outside medicine, Think-on-Graph~\cite{sun2024thinkongraph} lets an LLM
explore a knowledge graph step by step with beam search, and Microsoft's
GraphRAG~\cite{edge2024local} builds a graph of entities and communities from
text to answer global questions about a document collection.

All of these approaches use the knowledge graph only at \emph{inference
time} and assume that the facts needed to answer a question are present in
the graph. Zhou et al.~\cite{zhou2026breaks} questioned this assumption with
the BRINK benchmark (Section~\ref{sec:bg-rag}): when the triple that directly
supports an answer is removed, current KG-RAG methods show limited reasoning
ability, often fall back on knowledge memorised during pre-training, and
generalise differently depending on their design. This finding is highly
relevant for medical knowledge graphs, which are incomplete by nature
(Section~\ref{sec:bg-owa}).

\begin{bgexample}[Retrieval versus learnt reasoning under missing knowledge]
Assume the guideline knowledge graph contains the facts
\emph{(Drug~X, belongsToClass, MRA)} and
\emph{(MRA, recommendedFor, HFrEF)}, but \emph{not} the direct fact
\emph{(Drug~X, recommendedFor, HFrEF)}. A patient asks:
\emph{``Is Drug~X recommended for my type of heart failure?''}
\begin{itemize}
  \item A \textbf{retrieval-only} system searches for the direct triple, does
        not find it, and either answers ``no information'' or falls back on
        what the LLM memorised during pre-training.
  \item A model that was \textbf{fine-tuned on rule-based reasoning chains},
        e.g.\ derived from the rule
        $\mathit{belongsToClass}(x,c) \wedge \mathit{recommendedFor}(c,d)
        \Rightarrow \mathit{recommendedFor}(x,d)$,
        has learnt to combine the two remaining facts and can answer:
        \emph{``Drug~X is an MRA, and MRAs are recommended for HFrEF;
        therefore \ldots''}, while showing the facts it used.
\end{itemize}
This illustrates why this thesis combines knowledge-graph context at
inference time with reasoning patterns learnt from symbolic rules.
\end{bgexample}

\begin{rwgap}
Biomedical KG-enhanced LLMs ground answers in retrieved facts but do not
teach the model how to \emph{reason} over the graph when facts are missing.
Moreover, they target clinicians or benchmark questions, not patient
education based on clinical practice guidelines and patient questionnaires.
\end{rwgap}

% ---------------------------------------------------------------------
\section{Symbolic Rules and Neuro-Symbolic Reasoning with LLMs}
\label{sec:rw-nesy}
% ---------------------------------------------------------------------

\paragraph{Rule learning over knowledge graphs.}
Learning logical rules from knowledge graphs has a long tradition. The AMIE
family~\cite{galarraga2013amie,galarraga2015fast,lajus2020fast} mines closed
Horn rules under the Partial Completeness Assumption and is the rule miner
used in this thesis (Section~\ref{sec:bg-rule-mining}). AnyBURL~\cite{meilicke2019anytime}
learns rules bottom-up by sampling paths from the graph and generalising them,
which makes it efficient for knowledge-graph completion. Neural approaches
such as Neural~LP~\cite{yang2017differentiable} learn rules in a
differentiable way, trading some interpretability for end-to-end training.
In all of these works, rules are used to \emph{complete} the knowledge graph,
not to train a language model.

\paragraph{LLMs and rules.}
More recent work connects rules with LLMs. ChatRule~\cite{luo2025chatrule}
uses an LLM to \emph{generate} candidate rules from the semantic and
structural information of a knowledge graph and ranks them with
measures computed on the existing facts. Reasoning on Graphs
(RoG)~\cite{luo2024reasoning} fine-tunes an LLM to first generate relation
paths that are grounded in the knowledge graph (``planning'') and then to
reason over the retrieved paths, producing faithful and interpretable
answers. Closest to this thesis is Chain-of-Knowledge~\cite{zhang2024chain},
which mines rules from knowledge graphs, turns them into a knowledge-reasoning
dataset (KnowReason), and fine-tunes LLMs on it; to prevent the model from
simply memorising the rules, the authors add a trial-and-error mechanism.

\paragraph{Reasoning chains as training data.}
Training models on explicit reasoning chains has been shown to be
effective in general. Hsieh et al.~\cite{hsieh2023distilling} demonstrated
that smaller models fine-tuned on rationales extracted from a large LLM can
outperform much larger models while using less training data. Rule
groundings from a knowledge graph provide such rationales without a teacher
LLM, and their correctness can be traced back to explicit facts
(Section~\ref{sec:bg-rules-to-cot}).

\begin{rwgap}
Neuro-symbolic approaches that combine rules and LLMs have so far been
evaluated on open-domain benchmark knowledge graphs (e.g.\ Freebase-based
question-answering datasets).
It is unclear whether rule-guided chain-of-thought fine-tuning improves LLM
performance in the medical domain, where knowledge comes from heterogeneous
sources (guidelines, terminologies and patient questionnaires) and where
rules must be filtered by confidence before they are used.
\end{rwgap}

% ---------------------------------------------------------------------
\section{Automatic Prompt Optimisation}
\label{sec:rw-prompt}
% ---------------------------------------------------------------------

Since the quality of LLM outputs depends strongly on the prompt, several
methods search for better prompts automatically. Automatic Prompt Engineer
(APE)~\cite{zhou2023large} lets an LLM propose candidate instructions from
input--output examples and selects the best one by its score. ProTeGi
(APO)~\cite{pryzant2023automatic} improves prompts with natural-language
``gradients'', and OPRO~\cite{yang2024large} uses an LLM as an optimiser that
sees the history of previous prompts and their scores. Evolutionary
approaches maintain a population of prompts: EvoPrompt~\cite{guo2024connecting}
uses an LLM to perform mutation and crossover and improved over
human-written prompts on 31 datasets, and Promptbreeder~\cite{fernando2024promptbreeder}
additionally evolves the mutation prompts themselves. GAAPO~\cite{secheresse2025gaapo}
combines mutation, crossover, few-shot augmentation, APO and OPRO in a single
genetic algorithm and records which generator produced each improvement
(Section~\ref{sec:bg-apo}).

In medicine, Medprompt~\cite{nori2023can} combined dynamic few-shot
selection, self-generated chain-of-thought explanations and choice-shuffle
ensembling, which allowed GPT-4 to reach 90.2\,\% accuracy on the MedQA
benchmark without any domain-specific fine-tuning. Medprompt is a fixed
composition of prompting techniques, not a search over prompts, and like most
prompt-optimisation methods it was evaluated on multiple-choice questions.

\begin{rwgap}
Prompt-optimisation methods are mostly evaluated on classification and
multiple-choice tasks with very large models. They have not been studied for
generating \emph{structured} patient-education content, such as video
storyboards, with a small, domain-fine-tuned model whose medical content must
not be changed by the optimised prompt.
\end{rwgap}

% ---------------------------------------------------------------------
\section{Automatic Generation of Educational Videos}
\label{sec:rw-video}
% ---------------------------------------------------------------------

Video is an effective format for patient education: a systematic review
found strong evidence that video-based interventions improve the knowledge
of patients with chronic illnesses~\cite{deshpande2023video}. Producing such
videos, however, requires medical experts, script writers and designers.
Recent work therefore investigates how LLMs can automate parts of this
process. TheoremExplainAgent~\cite{ku2025theoremexplainagent} uses LLM agents
to plan and write Manim animation code for explanation videos of
mathematical and scientific theorems; the best configuration produced
videos for 93.8\,\% of 240 theorems, although most videos still contained
minor visual layout problems. Code2Video~\cite{chen2026code2video} also
generates educational videos as executable code, using a planner, a coder and
a critic agent, and reports a 40\,\% improvement over direct code generation.
In healthcare, Haider et al.~\cite{haider2026ai} compared AI-generated avatar
videos with text handouts for postoperative education in a pilot randomised
trial. Knowledge scores did not differ significantly, but participants rated
the videos as clearer and more engaging.

These systems focus on the \emph{production} of a video from a given topic
and rely on the general knowledge of the LLM for its content. They do not
verify the medical statements against a trusted knowledge source, and they do
not adapt the content to the needs of a specific patient population.

\begin{rwgap}
LLM-based video generation has been studied mainly for STEM topics, while
medical studies use generic AI tools without knowledge grounding. A method is
missing that produces patient-oriented video storyboards whose medical
content is derived from guideline-based knowledge and whose structure can be
validated before production.
\end{rwgap}

% ---------------------------------------------------------------------
\section{Summary and Research Gap}
\label{sec:rw-summary}
% ---------------------------------------------------------------------

Table~\ref{tab:rw-comparison} compares representative approaches along the
dimensions that are relevant for this thesis. Each line of research covers
only part of the problem: medical LLMs and patient-education studies provide
patient-oriented language but no explicit knowledge; biomedical KG-RAG
systems ground answers in knowledge graphs but only at inference time;
neuro-symbolic methods learn from rules but have not been applied to
medicine; prompt optimisation and video generation have not been combined
with knowledge-grounded medical content.

\begin{table}[tbp]
  \centering
  \small
  \setlength{\tabcolsep}{4pt}
  \begin{tabular}{@{}lccccccc@{}}
    \toprule
    \textbf{Approach} & \textbf{Med.} & \textbf{KG} & \textbf{Rules}
      & \textbf{CoT-FT} & \textbf{P-Opt} & \textbf{Patient} & \textbf{Video} \\
    \midrule
    Med-PaLM~\cite{singhal2023large}                 & \rwyes & \rwno  & \rwno  & \rwno  & \rwno  & \rwpart & \rwno \\
    King et al.~\cite{king2025improving}             & \rwyes & \rwno  & \rwno  & \rwno  & \rwno  & \rwyes  & \rwno \\
    MedRAG~\cite{xiong2024benchmarking}              & \rwyes & \rwno  & \rwno  & \rwno  & \rwno  & \rwno   & \rwno \\
    KG-RAG~\cite{soman2024biomedical}                & \rwyes & \rwyes & \rwno  & \rwno  & \rwno  & \rwno   & \rwno \\
    DR.KNOWS~\cite{gao2025leveraging}                & \rwyes & \rwyes & \rwno  & \rwno  & \rwno  & \rwno   & \rwno \\
    MedGraphRAG~\cite{wu2025medical}                 & \rwyes & \rwyes & \rwno  & \rwno  & \rwno  & \rwno   & \rwno \\
    RoG~\cite{luo2024reasoning}                      & \rwno  & \rwyes & \rwpart& \rwpart& \rwno  & \rwno   & \rwno \\
    ChatRule~\cite{luo2025chatrule}                  & \rwno  & \rwyes & \rwyes & \rwno  & \rwno  & \rwno   & \rwno \\
    Chain-of-Knowledge~\cite{zhang2024chain}         & \rwno  & \rwyes & \rwyes & \rwyes & \rwno  & \rwno   & \rwno \\
    Medprompt~\cite{nori2023can}                     & \rwyes & \rwno  & \rwno  & \rwno  & \rwpart& \rwno   & \rwno \\
    GAAPO~\cite{secheresse2025gaapo}                 & \rwno  & \rwno  & \rwno  & \rwno  & \rwyes & \rwno   & \rwno \\
    TheoremExplainAgent~\cite{ku2025theoremexplainagent} & \rwno & \rwno & \rwno & \rwno & \rwno & \rwno   & \rwyes \\
    Code2Video~\cite{chen2026code2video}             & \rwno  & \rwno  & \rwno  & \rwno  & \rwno  & \rwno   & \rwyes \\
    \midrule
    \textbf{This thesis}                             & \rwyes & \rwyes & \rwyes & \rwyes & \rwyes & \rwyes  & \rwyes \\
    \bottomrule
  \end{tabular}
  \caption[Comparison of related approaches]{Comparison of related
  approaches. \textbf{Med.}: medical domain; \textbf{KG}: uses a knowledge
  graph; \textbf{Rules}: uses symbolic rules; \textbf{CoT-FT}: fine-tunes the
  LLM on (rule-based) reasoning chains; \textbf{P-Opt}: automatic prompt
  optimisation; \textbf{Patient}: patient-oriented output; \textbf{Video}:
  educational video or storyboard generation. \rwyes~=~yes,
  \rwpart~=~partially, \rwno~=~no. For this thesis, \textbf{Video} refers to
  storyboard scripts as input for video production.}
  \label{tab:rw-comparison}
\end{table}

Based on this analysis, the research gap addressed in this thesis can be
summarised as follows.

\begin{rwgap}[Research Gap of this Thesis]
There is no approach that (i)~represents clinical guidelines and patient
questionnaires about cardiovascular disease as knowledge graphs,
(ii)~turns symbolic rules mined from these graphs into chain-of-thought
training data for a medical LLM, (iii)~evaluates whether this improves
reasoning when knowledge is incomplete, and (iv)~uses the resulting model,
together with an automatically optimised prompt, to generate
patient-oriented storyboards for educational videos. This thesis addresses
this gap with the semantic-driven approach presented in
Chapter~4.
\end{rwgap}

%  NEEDS:  bg_preamble.tex, rw_preamble.tex, bg_references.bib,
%          rw_references.bib
% =====================================================================

\chapter{Related Work}
\label{chap:related-work}

This chapter reviews the research that this thesis builds on and
distinguishes it from existing approaches. The thesis lies at the
intersection of three research areas (Figure~\ref{fig:rw-positioning}):
(i)~large language models for medicine and patient education,
(ii)~knowledge graphs and neuro-symbolic reasoning for LLMs, and
(iii)~automatic prompt optimisation and the automatic generation of
educational videos. Section~\ref{sec:rw-medical-llm} discusses medical LLMs
and their limitations for patient education.
Section~\ref{sec:rw-kg-llm} reviews approaches that ground medical LLMs in
knowledge graphs, and Section~\ref{sec:rw-nesy} reviews methods that combine
symbolic rules with LLM reasoning. Section~\ref{sec:rw-prompt} covers
automatic prompt optimisation, and Section~\ref{sec:rw-video} covers the
generation of educational videos. Section~\ref{sec:rw-summary} compares all
approaches and summarises the research gap addressed by this thesis.

\begin{figure}[htbp]
  \centering
  \begin{tikzpicture}[font=\footnotesize]
    \def\r{3.0}
    \begin{scope}[blend group=multiply]
      \fill[bgOrangeLight] (0,0) circle (\r);
      \fill[bgBlueLight]   (3.4,0) circle (\r);
      \fill[bgPurpleLight] (1.7,-3.0) circle (\r);
    \end{scope}
    \draw[bgOrange, thick] (0,0) circle (\r);
    \draw[bgBlue, thick]   (3.4,0) circle (\r);
    \draw[bgPurple, thick] (1.7,-3.0) circle (\r);
    % area titles (outside the circles)
    \node[align=center, text=bgOrange, font=\footnotesize\bfseries] at (-1.4,3.45) {LLMs for medicine\\\& patient education};
    \node[align=center, text=bgBlue, font=\footnotesize\bfseries]   at (4.8,3.45) {Knowledge graphs \&\\neuro-symbolic AI};
    \node[align=center, text=bgPurple, font=\footnotesize\bfseries] at (1.7,-6.5) {Prompt optimisation \&\\educational video generation};
    % works that belong to one area only
    \node[align=center, font=\scriptsize] at (-1.25,0.1) {Med-PaLM~\cite{singhal2023large}\\Med-HALT~\cite{pal2023medhalt}\\King et al.~\cite{king2025improving}};
    \node[align=center, font=\scriptsize] at (4.65,0.1) {AMIE~3~\cite{lajus2020fast}\\RoG~\cite{luo2024reasoning}\\ChatRule~\cite{luo2025chatrule}\\BRINK~\cite{zhou2026breaks}};
    \node[align=center, font=\scriptsize] at (1.7,-4.35) {GAAPO~\cite{secheresse2025gaapo}\\OPRO~\cite{yang2024large}\\Code2Video~\cite{chen2026code2video}};
    % overlaps
    \node[align=center, font=\tiny] at (1.7,0.85) {KG-RAG~\cite{soman2024biomedical}\\MedGraphRAG~\cite{wu2025medical}\\DR.KNOWS~\cite{gao2025leveraging}};
    \node[align=center, font=\tiny] at (-0.1,-2.05) {Medprompt\\\cite{nori2023can}};
    % centre
    \node[draw=black, thick, fill=white, rounded corners=3pt, align=center,
          font=\scriptsize\bfseries, inner sep=3pt] at (1.7,-0.95) {This\\thesis};
  \end{tikzpicture}
  \caption[Positioning of this thesis within related research areas]{Positioning
  of this thesis. Existing work covers one or two of the three research areas;
  this thesis combines guideline-based knowledge graphs, rule-guided
  reasoning, prompt optimisation and patient-oriented storyboards for
  educational videos.}
  \label{fig:rw-positioning}
\end{figure}

% ---------------------------------------------------------------------
\section{Large Language Models in Medicine and Patient Education}
\label{sec:rw-medical-llm}
% ---------------------------------------------------------------------

\paragraph{Medical question answering.}
Singhal et al.~\cite{singhal2023large} introduced MultiMedQA, a benchmark that
combines professional medical exams, research questions and consumer health
questions, and adapted the general-purpose model Flan-PaLM to the medical
domain through instruction prompt tuning. In a human evaluation, 92.6\,\% of
the long-form answers of the resulting model, Med-PaLM, were judged to be in
line with scientific consensus, compared with 92.9\,\% of clinician answers;
nevertheless, the model remained inferior to clinicians on several other
evaluation axes. Its successor
Med-PaLM~2~\cite{singhal2025toward} further narrowed this gap. Nori et
al.~\cite{nori2023can} showed that a general model (GPT-4) can reach similar
results without any medical fine-tuning when it is combined with a carefully
designed prompting strategy (Section~\ref{sec:rw-prompt}). Surveys of the
field~\cite{thirunavukarasu2023large} conclude that LLMs are promising for
clinical documentation, education and patient communication, but that their
reliability, transparency and accountability are not yet sufficient for
unsupervised use.

\paragraph{Patient communication and patient education.}
LLMs have also been evaluated directly on patient-facing tasks. Ayers et
al.~\cite{ayers2023comparing} compared answers by physicians and by ChatGPT to
195 patient questions from a public online forum; in 78.6\,\% of 585
evaluations, healthcare professionals preferred the chatbot answer, which was
rated higher for both quality and empathy. In cardiology, King et
al.~\cite{king2025improving} used GPT-4 to simplify 143 heart-failure patient
education materials from leading US cardiology institutions. The median
reading level decreased from grade~10.3 to grade~7.3, and a cardiologist rated
all simplified texts as equally accurate as the originals. These studies show
that LLMs can make medical information more accessible. However, they rely
on general-purpose models without explicit medical knowledge, they evaluate
plain text only, and correctness is checked manually after generation.

\paragraph{Hallucinations in medical LLMs.}
The main obstacle to the safe use of LLMs for patients is hallucination
(Section~\ref{sec:bg-hallucination}). Pal et al.~\cite{pal2023medhalt}
proposed Med-HALT, a benchmark that tests both reasoning hallucinations (e.g.\
choosing an answer to a question with a false premise) and memory-based
hallucinations (e.g.\ inventing references), and found that all evaluated
LLMs were affected. Kim et al.~\cite{kim2025medicalhallucinationsfoundationmodels}
analysed eleven foundation models on seven medical tasks and, together with a
survey of clinicians, concluded that medical hallucinations are frequent in
practice. Interestingly, they also observed that chain-of-thought prompting
reduced hallucinations in most comparisons, which supports the idea of making
the reasoning of a model explicit.

\begin{rwgap}
Medical LLMs can produce fluent and empathetic answers, but they are not
explicitly grounded in clinical guidelines, their reasoning is not
verifiable, and existing patient-education studies focus on text.
Approaches are needed that connect patient-oriented generation with
explicit, guideline-based medical knowledge.
\end{rwgap}

% ---------------------------------------------------------------------
\section{Knowledge-Graph-Enhanced LLMs in the Biomedical Domain}
\label{sec:rw-kg-llm}
% ---------------------------------------------------------------------

A common way to reduce hallucinations is to supply LLMs with external
knowledge at inference time (Section~\ref{sec:bg-rag}). Xiong et
al.~\cite{xiong2024benchmarking} systematically benchmarked
retrieval-augmented generation for medicine on MIRAGE, a collection of 7,663
questions from five medical QA datasets. Their MedRAG toolkit improved the
accuracy of six LLMs by up to 18\,\% compared with chain-of-thought prompting.
These systems retrieve unstructured text passages, for example from PubMed or
textbooks.

Several works retrieve structured knowledge from \emph{knowledge graphs}
instead. Soman et al.~\cite{soman2024biomedical} proposed KG-RAG, which
extracts disease-specific context from SPOKE, a large biomedical knowledge
graph, and adds it to the prompt. By pruning the retrieved context with
embeddings, they reduced the number of prompt tokens by more than 50\,\%
while improving the answers of models such as Llama-2 and GPT-4. Gao et
al.~\cite{gao2025leveraging} presented DR.KNOWS, which selects paths from a
UMLS-based knowledge graph that are relevant to a patient's electronic health
record and uses them to support diagnosis prediction; clinicians judged the
resulting diagnostic rationales to be clinically sound. Wu et
al.~\cite{wu2025medical} proposed MedGraphRAG, which links private medical
documents to trusted medical sources and vocabularies in a three-level graph,
and retrieves from this graph so that answers are supported by evidence.
Outside medicine, Think-on-Graph~\cite{sun2024thinkongraph} lets an LLM
explore a knowledge graph step by step with beam search, and Microsoft's
GraphRAG~\cite{edge2024local} builds a graph of entities and communities from
text to answer global questions about a document collection.

All of these approaches use the knowledge graph only at \emph{inference
time} and assume that the facts needed to answer a question are present in
the graph. Zhou et al.~\cite{zhou2026breaks} questioned this assumption with
the BRINK benchmark (Section~\ref{sec:bg-rag}): when the triple that directly
supports an answer is removed, current KG-RAG methods show limited reasoning
ability, often fall back on knowledge memorised during pre-training, and
generalise differently depending on their design. This finding is highly
relevant for medical knowledge graphs, which are incomplete by nature
(Section~\ref{sec:bg-owa}).

\begin{bgexample}[Retrieval versus learnt reasoning under missing knowledge]
Assume the guideline knowledge graph contains the facts
\emph{(Drug~X, belongsToClass, MRA)} and
\emph{(MRA, recommendedFor, HFrEF)}, but \emph{not} the direct fact
\emph{(Drug~X, recommendedFor, HFrEF)}. A patient asks:
\emph{``Is Drug~X recommended for my type of heart failure?''}
\begin{itemize}
  \item A \textbf{retrieval-only} system searches for the direct triple, does
        not find it, and either answers ``no information'' or falls back on
        what the LLM memorised during pre-training.
  \item A model that was \textbf{fine-tuned on rule-based reasoning chains},
        e.g.\ derived from the rule
        $\mathit{belongsToClass}(x,c) \wedge \mathit{recommendedFor}(c,d)
        \Rightarrow \mathit{recommendedFor}(x,d)$,
        has learnt to combine the two remaining facts and can answer:
        \emph{``Drug~X is an MRA, and MRAs are recommended for HFrEF;
        therefore \ldots''}, while showing the facts it used.
\end{itemize}
This illustrates why this thesis combines knowledge-graph context at
inference time with reasoning patterns learnt from symbolic rules.
\end{bgexample}

\begin{rwgap}
Biomedical KG-enhanced LLMs ground answers in retrieved facts but do not
teach the model how to \emph{reason} over the graph when facts are missing.
Moreover, they target clinicians or benchmark questions, not patient
education based on clinical practice guidelines and patient questionnaires.
\end{rwgap}

% ---------------------------------------------------------------------
\section{Symbolic Rules and Neuro-Symbolic Reasoning with LLMs}
\label{sec:rw-nesy}
% ---------------------------------------------------------------------

\paragraph{Rule learning over knowledge graphs.}
Learning logical rules from knowledge graphs has a long tradition. The AMIE
family~\cite{galarraga2013amie,galarraga2015fast,lajus2020fast} mines closed
Horn rules under the Partial Completeness Assumption and is the rule miner
used in this thesis (Section~\ref{sec:bg-rule-mining}). AnyBURL~\cite{meilicke2019anytime}
learns rules bottom-up by sampling paths from the graph and generalising them,
which makes it efficient for knowledge-graph completion. Neural approaches
such as Neural~LP~\cite{yang2017differentiable} learn rules in a
differentiable way, trading some interpretability for end-to-end training.
In all of these works, rules are used to \emph{complete} the knowledge graph,
not to train a language model.

\paragraph{LLMs and rules.}
More recent work connects rules with LLMs. ChatRule~\cite{luo2025chatrule}
uses an LLM to \emph{generate} candidate rules from the semantic and
structural information of a knowledge graph and ranks them with
measures computed on the existing facts. Reasoning on Graphs
(RoG)~\cite{luo2024reasoning} fine-tunes an LLM to first generate relation
paths that are grounded in the knowledge graph (``planning'') and then to
reason over the retrieved paths, producing faithful and interpretable
answers. Closest to this thesis is Chain-of-Knowledge~\cite{zhang2024chain},
which mines rules from knowledge graphs, turns them into a knowledge-reasoning
dataset (KnowReason), and fine-tunes LLMs on it; to prevent the model from
simply memorising the rules, the authors add a trial-and-error mechanism.

\paragraph{Reasoning chains as training data.}
Training models on explicit reasoning chains has been shown to be
effective in general. Hsieh et al.~\cite{hsieh2023distilling} demonstrated
that smaller models fine-tuned on rationales extracted from a large LLM can
outperform much larger models while using less training data. Rule
groundings from a knowledge graph provide such rationales without a teacher
LLM, and their correctness can be traced back to explicit facts
(Section~\ref{sec:bg-rules-to-cot}).

\begin{rwgap}
Neuro-symbolic approaches that combine rules and LLMs have so far been
evaluated on open-domain benchmark knowledge graphs (e.g.\ Freebase-based
question-answering datasets).
It is unclear whether rule-guided chain-of-thought fine-tuning improves LLM
performance in the medical domain, where knowledge comes from heterogeneous
sources (guidelines, terminologies and patient questionnaires) and where
rules must be filtered by confidence before they are used.
\end{rwgap}

% ---------------------------------------------------------------------
\section{Automatic Prompt Optimisation}
\label{sec:rw-prompt}
% ---------------------------------------------------------------------

Since the quality of LLM outputs depends strongly on the prompt, several
methods search for better prompts automatically. Automatic Prompt Engineer
(APE)~\cite{zhou2023large} lets an LLM propose candidate instructions from
input--output examples and selects the best one by its score. ProTeGi
(APO)~\cite{pryzant2023automatic} improves prompts with natural-language
``gradients'', and OPRO~\cite{yang2024large} uses an LLM as an optimiser that
sees the history of previous prompts and their scores. Evolutionary
approaches maintain a population of prompts: EvoPrompt~\cite{guo2024connecting}
uses an LLM to perform mutation and crossover and improved over
human-written prompts on 31 datasets, and Promptbreeder~\cite{fernando2024promptbreeder}
additionally evolves the mutation prompts themselves. GAAPO~\cite{secheresse2025gaapo}
combines mutation, crossover, few-shot augmentation, APO and OPRO in a single
genetic algorithm and records which generator produced each improvement
(Section~\ref{sec:bg-apo}).

In medicine, Medprompt~\cite{nori2023can} combined dynamic few-shot
selection, self-generated chain-of-thought explanations and choice-shuffle
ensembling, which allowed GPT-4 to reach 90.2\,\% accuracy on the MedQA
benchmark without any domain-specific fine-tuning. Medprompt is a fixed
composition of prompting techniques, not a search over prompts, and like most
prompt-optimisation methods it was evaluated on multiple-choice questions.

\begin{rwgap}
Prompt-optimisation methods are mostly evaluated on classification and
multiple-choice tasks with very large models. They have not been studied for
generating \emph{structured} patient-education content, such as video
storyboards, with a small, domain-fine-tuned model whose medical content must
not be changed by the optimised prompt.
\end{rwgap}

% ---------------------------------------------------------------------
\section{Automatic Generation of Educational Videos}
\label{sec:rw-video}
% ---------------------------------------------------------------------

Video is an effective format for patient education: a systematic review
found strong evidence that video-based interventions improve the knowledge
of patients with chronic illnesses~\cite{deshpande2023video}. Producing such
videos, however, requires medical experts, script writers and designers.
Recent work therefore investigates how LLMs can automate parts of this
process. TheoremExplainAgent~\cite{ku2025theoremexplainagent} uses LLM agents
to plan and write Manim animation code for explanation videos of
mathematical and scientific theorems; the best configuration produced
videos for 93.8\,\% of 240 theorems, although most videos still contained
minor visual layout problems. Code2Video~\cite{chen2026code2video} also
generates educational videos as executable code, using a planner, a coder and
a critic agent, and reports a 40\,\% improvement over direct code generation.
In healthcare, Haider et al.~\cite{haider2026ai} compared AI-generated avatar
videos with text handouts for postoperative education in a pilot randomised
trial. Knowledge scores did not differ significantly, but participants rated
the videos as clearer and more engaging.

These systems focus on the \emph{production} of a video from a given topic
and rely on the general knowledge of the LLM for its content. They do not
verify the medical statements against a trusted knowledge source, and they do
not adapt the content to the needs of a specific patient population.

\begin{rwgap}
LLM-based video generation has been studied mainly for STEM topics, while
medical studies use generic AI tools without knowledge grounding. A method is
missing that produces patient-oriented video storyboards whose medical
content is derived from guideline-based knowledge and whose structure can be
validated before production.
\end{rwgap}

% ---------------------------------------------------------------------
\section{Summary and Research Gap}
\label{sec:rw-summary}
% ---------------------------------------------------------------------

Table~\ref{tab:rw-comparison} compares representative approaches along the
dimensions that are relevant for this thesis. Each line of research covers
only part of the problem: medical LLMs and patient-education studies provide
patient-oriented language but no explicit knowledge; biomedical KG-RAG
systems ground answers in knowledge graphs but only at inference time;
neuro-symbolic methods learn from rules but have not been applied to
medicine; prompt optimisation and video generation have not been combined
with knowledge-grounded medical content.

\begin{table}[tbp]
  \centering
  \small
  \setlength{\tabcolsep}{4pt}
  \begin{tabular}{@{}lccccccc@{}}
    \toprule
    \textbf{Approach} & \textbf{Med.} & \textbf{KG} & \textbf{Rules}
      & \textbf{CoT-FT} & \textbf{P-Opt} & \textbf{Patient} & \textbf{Video} \\
    \midrule
    Med-PaLM~\cite{singhal2023large}                 & \rwyes & \rwno  & \rwno  & \rwno  & \rwno  & \rwpart & \rwno \\
    King et al.~\cite{king2025improving}             & \rwyes & \rwno  & \rwno  & \rwno  & \rwno  & \rwyes  & \rwno \\
    MedRAG~\cite{xiong2024benchmarking}              & \rwyes & \rwno  & \rwno  & \rwno  & \rwno  & \rwno   & \rwno \\
    KG-RAG~\cite{soman2024biomedical}                & \rwyes & \rwyes & \rwno  & \rwno  & \rwno  & \rwno   & \rwno \\
    DR.KNOWS~\cite{gao2025leveraging}                & \rwyes & \rwyes & \rwno  & \rwno  & \rwno  & \rwno   & \rwno \\
    MedGraphRAG~\cite{wu2025medical}                 & \rwyes & \rwyes & \rwno  & \rwno  & \rwno  & \rwno   & \rwno \\
    RoG~\cite{luo2024reasoning}                      & \rwno  & \rwyes & \rwpart& \rwpart& \rwno  & \rwno   & \rwno \\
    ChatRule~\cite{luo2025chatrule}                  & \rwno  & \rwyes & \rwyes & \rwno  & \rwno  & \rwno   & \rwno \\
    Chain-of-Knowledge~\cite{zhang2024chain}         & \rwno  & \rwyes & \rwyes & \rwyes & \rwno  & \rwno   & \rwno \\
    Medprompt~\cite{nori2023can}                     & \rwyes & \rwno  & \rwno  & \rwno  & \rwpart& \rwno   & \rwno \\
    GAAPO~\cite{secheresse2025gaapo}                 & \rwno  & \rwno  & \rwno  & \rwno  & \rwyes & \rwno   & \rwno \\
    TheoremExplainAgent~\cite{ku2025theoremexplainagent} & \rwno & \rwno & \rwno & \rwno & \rwno & \rwno   & \rwyes \\
    Code2Video~\cite{chen2026code2video}             & \rwno  & \rwno  & \rwno  & \rwno  & \rwno  & \rwno   & \rwyes \\
    \midrule
    \textbf{This thesis}                             & \rwyes & \rwyes & \rwyes & \rwyes & \rwyes & \rwyes  & \rwyes \\
    \bottomrule
  \end{tabular}
  \caption[Comparison of related approaches]{Comparison of related
  approaches. \textbf{Med.}: medical domain; \textbf{KG}: uses a knowledge
  graph; \textbf{Rules}: uses symbolic rules; \textbf{CoT-FT}: fine-tunes the
  LLM on (rule-based) reasoning chains; \textbf{P-Opt}: automatic prompt
  optimisation; \textbf{Patient}: patient-oriented output; \textbf{Video}:
  educational video or storyboard generation. \rwyes~=~yes,
  \rwpart~=~partially, \rwno~=~no. For this thesis, \textbf{Video} refers to
  storyboard scripts as input for video production.}
  \label{tab:rw-comparison}
\end{table}

Based on this analysis, the research gap addressed in this thesis can be
summarised as follows.

\begin{rwgap}[Research Gap of this Thesis]
There is no approach that (i)~represents clinical guidelines and patient
questionnaires about cardiovascular disease as knowledge graphs,
(ii)~turns symbolic rules mined from these graphs into chain-of-thought
training data for a medical LLM, (iii)~evaluates whether this improves
reasoning when knowledge is incomplete, and (iv)~uses the resulting model,
together with an automatically optimised prompt, to generate
patient-oriented storyboards for educational videos. This thesis addresses
this gap with the semantic-driven approach presented in
Chapter~4.
\end{rwgap}

%  NEEDS:  bg_preamble.tex, rw_preamble.tex, bg_references.bib,
%          rw_references.bib
% =====================================================================

\chapter{Related Work}
\label{chap:related-work}

This chapter reviews the research that this thesis builds on and
distinguishes it from existing approaches. The thesis lies at the
intersection of three research areas (Figure~\ref{fig:rw-positioning}):
(i)~large language models for medicine and patient education,
(ii)~knowledge graphs and neuro-symbolic reasoning for LLMs, and
(iii)~automatic prompt optimisation and the automatic generation of
educational videos. Section~\ref{sec:rw-medical-llm} discusses medical LLMs
and their limitations for patient education.
Section~\ref{sec:rw-kg-llm} reviews approaches that ground medical LLMs in
knowledge graphs, and Section~\ref{sec:rw-nesy} reviews methods that combine
symbolic rules with LLM reasoning. Section~\ref{sec:rw-prompt} covers
automatic prompt optimisation, and Section~\ref{sec:rw-video} covers the
generation of educational videos. Section~\ref{sec:rw-summary} compares all
approaches and summarises the research gap addressed by this thesis.

\begin{figure}[htbp]
  \centering
  \begin{tikzpicture}[font=\footnotesize]
    \def\r{3.0}
    \begin{scope}[blend group=multiply]
      \fill[bgOrangeLight] (0,0) circle (\r);
      \fill[bgBlueLight]   (3.4,0) circle (\r);
      \fill[bgPurpleLight] (1.7,-3.0) circle (\r);
    \end{scope}
    \draw[bgOrange, thick] (0,0) circle (\r);
    \draw[bgBlue, thick]   (3.4,0) circle (\r);
    \draw[bgPurple, thick] (1.7,-3.0) circle (\r);
    % area titles (outside the circles)
    \node[align=center, text=bgOrange, font=\footnotesize\bfseries] at (-1.4,3.45) {LLMs for medicine\\\& patient education};
    \node[align=center, text=bgBlue, font=\footnotesize\bfseries]   at (4.8,3.45) {Knowledge graphs \&\\neuro-symbolic AI};
    \node[align=center, text=bgPurple, font=\footnotesize\bfseries] at (1.7,-6.5) {Prompt optimisation \&\\educational video generation};
    % works that belong to one area only
    \node[align=center, font=\scriptsize] at (-1.25,0.1) {Med-PaLM~\cite{singhal2023large}\\Med-HALT~\cite{pal2023medhalt}\\King et al.~\cite{king2025improving}};
    \node[align=center, font=\scriptsize] at (4.65,0.1) {AMIE~3~\cite{lajus2020fast}\\RoG~\cite{luo2024reasoning}\\ChatRule~\cite{luo2025chatrule}\\BRINK~\cite{zhou2026breaks}};
    \node[align=center, font=\scriptsize] at (1.7,-4.35) {GAAPO~\cite{secheresse2025gaapo}\\OPRO~\cite{yang2024large}\\Code2Video~\cite{chen2026code2video}};
    % overlaps
    \node[align=center, font=\tiny] at (1.7,0.85) {KG-RAG~\cite{soman2024biomedical}\\MedGraphRAG~\cite{wu2025medical}\\DR.KNOWS~\cite{gao2025leveraging}};
    \node[align=center, font=\tiny] at (-0.1,-2.05) {Medprompt\\\cite{nori2023can}};
    % centre
    \node[draw=black, thick, fill=white, rounded corners=3pt, align=center,
          font=\scriptsize\bfseries, inner sep=3pt] at (1.7,-0.95) {This\\thesis};
  \end{tikzpicture}
  \caption[Positioning of this thesis within related research areas]{Positioning
  of this thesis. Existing work covers one or two of the three research areas;
  this thesis combines guideline-based knowledge graphs, rule-guided
  reasoning, prompt optimisation and patient-oriented storyboards for
  educational videos.}
  \label{fig:rw-positioning}
\end{figure}

% ---------------------------------------------------------------------
\section{Large Language Models in Medicine and Patient Education}
\label{sec:rw-medical-llm}
% ---------------------------------------------------------------------

\paragraph{Medical question answering.}
Singhal et al.~\cite{singhal2023large} introduced MultiMedQA, a benchmark that
combines professional medical exams, research questions and consumer health
questions, and adapted the general-purpose model Flan-PaLM to the medical
domain through instruction prompt tuning. In a human evaluation, 92.6\,\% of
the long-form answers of the resulting model, Med-PaLM, were judged to be in
line with scientific consensus, compared with 92.9\,\% of clinician answers;
nevertheless, the model remained inferior to clinicians on several other
evaluation axes. Its successor
Med-PaLM~2~\cite{singhal2025toward} further narrowed this gap. Nori et
al.~\cite{nori2023can} showed that a general model (GPT-4) can reach similar
results without any medical fine-tuning when it is combined with a carefully
designed prompting strategy (Section~\ref{sec:rw-prompt}). Surveys of the
field~\cite{thirunavukarasu2023large} conclude that LLMs are promising for
clinical documentation, education and patient communication, but that their
reliability, transparency and accountability are not yet sufficient for
unsupervised use.

\paragraph{Patient communication and patient education.}
LLMs have also been evaluated directly on patient-facing tasks. Ayers et
al.~\cite{ayers2023comparing} compared answers by physicians and by ChatGPT to
195 patient questions from a public online forum; in 78.6\,\% of 585
evaluations, healthcare professionals preferred the chatbot answer, which was
rated higher for both quality and empathy. In cardiology, King et
al.~\cite{king2025improving} used GPT-4 to simplify 143 heart-failure patient
education materials from leading US cardiology institutions. The median
reading level decreased from grade~10.3 to grade~7.3, and a cardiologist rated
all simplified texts as equally accurate as the originals. These studies show
that LLMs can make medical information more accessible. However, they rely
on general-purpose models without explicit medical knowledge, they evaluate
plain text only, and correctness is checked manually after generation.

\paragraph{Hallucinations in medical LLMs.}
The main obstacle to the safe use of LLMs for patients is hallucination
(Section~\ref{sec:bg-hallucination}). Pal et al.~\cite{pal2023medhalt}
proposed Med-HALT, a benchmark that tests both reasoning hallucinations (e.g.\
choosing an answer to a question with a false premise) and memory-based
hallucinations (e.g.\ inventing references), and found that all evaluated
LLMs were affected. Kim et al.~\cite{kim2025medicalhallucinationsfoundationmodels}
analysed eleven foundation models on seven medical tasks and, together with a
survey of clinicians, concluded that medical hallucinations are frequent in
practice. Interestingly, they also observed that chain-of-thought prompting
reduced hallucinations in most comparisons, which supports the idea of making
the reasoning of a model explicit.

\begin{rwgap}
Medical LLMs can produce fluent and empathetic answers, but they are not
explicitly grounded in clinical guidelines, their reasoning is not
verifiable, and existing patient-education studies focus on text.
Approaches are needed that connect patient-oriented generation with
explicit, guideline-based medical knowledge.
\end{rwgap}

% ---------------------------------------------------------------------
\section{Knowledge-Graph-Enhanced LLMs in the Biomedical Domain}
\label{sec:rw-kg-llm}
% ---------------------------------------------------------------------

A common way to reduce hallucinations is to supply LLMs with external
knowledge at inference time (Section~\ref{sec:bg-rag}). Xiong et
al.~\cite{xiong2024benchmarking} systematically benchmarked
retrieval-augmented generation for medicine on MIRAGE, a collection of 7,663
questions from five medical QA datasets. Their MedRAG toolkit improved the
accuracy of six LLMs by up to 18\,\% compared with chain-of-thought prompting.
These systems retrieve unstructured text passages, for example from PubMed or
textbooks.

Several works retrieve structured knowledge from \emph{knowledge graphs}
instead. Soman et al.~\cite{soman2024biomedical} proposed KG-RAG, which
extracts disease-specific context from SPOKE, a large biomedical knowledge
graph, and adds it to the prompt. By pruning the retrieved context with
embeddings, they reduced the number of prompt tokens by more than 50\,\%
while improving the answers of models such as Llama-2 and GPT-4. Gao et
al.~\cite{gao2025leveraging} presented DR.KNOWS, which selects paths from a
UMLS-based knowledge graph that are relevant to a patient's electronic health
record and uses them to support diagnosis prediction; clinicians judged the
resulting diagnostic rationales to be clinically sound. Wu et
al.~\cite{wu2025medical} proposed MedGraphRAG, which links private medical
documents to trusted medical sources and vocabularies in a three-level graph,
and retrieves from this graph so that answers are supported by evidence.
Outside medicine, Think-on-Graph~\cite{sun2024thinkongraph} lets an LLM
explore a knowledge graph step by step with beam search, and Microsoft's
GraphRAG~\cite{edge2024local} builds a graph of entities and communities from
text to answer global questions about a document collection.

All of these approaches use the knowledge graph only at \emph{inference
time} and assume that the facts needed to answer a question are present in
the graph. Zhou et al.~\cite{zhou2026breaks} questioned this assumption with
the BRINK benchmark (Section~\ref{sec:bg-rag}): when the triple that directly
supports an answer is removed, current KG-RAG methods show limited reasoning
ability, often fall back on knowledge memorised during pre-training, and
generalise differently depending on their design. This finding is highly
relevant for medical knowledge graphs, which are incomplete by nature
(Section~\ref{sec:bg-owa}).

\begin{bgexample}[Retrieval versus learnt reasoning under missing knowledge]
Assume the guideline knowledge graph contains the facts
\emph{(Drug~X, belongsToClass, MRA)} and
\emph{(MRA, recommendedFor, HFrEF)}, but \emph{not} the direct fact
\emph{(Drug~X, recommendedFor, HFrEF)}. A patient asks:
\emph{``Is Drug~X recommended for my type of heart failure?''}
\begin{itemize}
  \item A \textbf{retrieval-only} system searches for the direct triple, does
        not find it, and either answers ``no information'' or falls back on
        what the LLM memorised during pre-training.
  \item A model that was \textbf{fine-tuned on rule-based reasoning chains},
        e.g.\ derived from the rule
        $\mathit{belongsToClass}(x,c) \wedge \mathit{recommendedFor}(c,d)
        \Rightarrow \mathit{recommendedFor}(x,d)$,
        has learnt to combine the two remaining facts and can answer:
        \emph{``Drug~X is an MRA, and MRAs are recommended for HFrEF;
        therefore \ldots''}, while showing the facts it used.
\end{itemize}
This illustrates why this thesis combines knowledge-graph context at
inference time with reasoning patterns learnt from symbolic rules.
\end{bgexample}

\begin{rwgap}
Biomedical KG-enhanced LLMs ground answers in retrieved facts but do not
teach the model how to \emph{reason} over the graph when facts are missing.
Moreover, they target clinicians or benchmark questions, not patient
education based on clinical practice guidelines and patient questionnaires.
\end{rwgap}

% ---------------------------------------------------------------------
\section{Symbolic Rules and Neuro-Symbolic Reasoning with LLMs}
\label{sec:rw-nesy}
% ---------------------------------------------------------------------

\paragraph{Rule learning over knowledge graphs.}
Learning logical rules from knowledge graphs has a long tradition. The AMIE
family~\cite{galarraga2013amie,galarraga2015fast,lajus2020fast} mines closed
Horn rules under the Partial Completeness Assumption and is the rule miner
used in this thesis (Section~\ref{sec:bg-rule-mining}). AnyBURL~\cite{meilicke2019anytime}
learns rules bottom-up by sampling paths from the graph and generalising them,
which makes it efficient for knowledge-graph completion. Neural approaches
such as Neural~LP~\cite{yang2017differentiable} learn rules in a
differentiable way, trading some interpretability for end-to-end training.
In all of these works, rules are used to \emph{complete} the knowledge graph,
not to train a language model.

\paragraph{LLMs and rules.}
More recent work connects rules with LLMs. ChatRule~\cite{luo2025chatrule}
uses an LLM to \emph{generate} candidate rules from the semantic and
structural information of a knowledge graph and ranks them with
measures computed on the existing facts. Reasoning on Graphs
(RoG)~\cite{luo2024reasoning} fine-tunes an LLM to first generate relation
paths that are grounded in the knowledge graph (``planning'') and then to
reason over the retrieved paths, producing faithful and interpretable
answers. Closest to this thesis is Chain-of-Knowledge~\cite{zhang2024chain},
which mines rules from knowledge graphs, turns them into a knowledge-reasoning
dataset (KnowReason), and fine-tunes LLMs on it; to prevent the model from
simply memorising the rules, the authors add a trial-and-error mechanism.

\paragraph{Reasoning chains as training data.}
Training models on explicit reasoning chains has been shown to be
effective in general. Hsieh et al.~\cite{hsieh2023distilling} demonstrated
that smaller models fine-tuned on rationales extracted from a large LLM can
outperform much larger models while using less training data. Rule
groundings from a knowledge graph provide such rationales without a teacher
LLM, and their correctness can be traced back to explicit facts
(Section~\ref{sec:bg-rules-to-cot}).

\begin{rwgap}
Neuro-symbolic approaches that combine rules and LLMs have so far been
evaluated on open-domain benchmark knowledge graphs (e.g.\ Freebase-based
question-answering datasets).
It is unclear whether rule-guided chain-of-thought fine-tuning improves LLM
performance in the medical domain, where knowledge comes from heterogeneous
sources (guidelines, terminologies and patient questionnaires) and where
rules must be filtered by confidence before they are used.
\end{rwgap}

% ---------------------------------------------------------------------
\section{Automatic Prompt Optimisation}
\label{sec:rw-prompt}
% ---------------------------------------------------------------------

Since the quality of LLM outputs depends strongly on the prompt, several
methods search for better prompts automatically. Automatic Prompt Engineer
(APE)~\cite{zhou2023large} lets an LLM propose candidate instructions from
input--output examples and selects the best one by its score. ProTeGi
(APO)~\cite{pryzant2023automatic} improves prompts with natural-language
``gradients'', and OPRO~\cite{yang2024large} uses an LLM as an optimiser that
sees the history of previous prompts and their scores. Evolutionary
approaches maintain a population of prompts: EvoPrompt~\cite{guo2024connecting}
uses an LLM to perform mutation and crossover and improved over
human-written prompts on 31 datasets, and Promptbreeder~\cite{fernando2024promptbreeder}
additionally evolves the mutation prompts themselves. GAAPO~\cite{secheresse2025gaapo}
combines mutation, crossover, few-shot augmentation, APO and OPRO in a single
genetic algorithm and records which generator produced each improvement
(Section~\ref{sec:bg-apo}).

In medicine, Medprompt~\cite{nori2023can} combined dynamic few-shot
selection, self-generated chain-of-thought explanations and choice-shuffle
ensembling, which allowed GPT-4 to reach 90.2\,\% accuracy on the MedQA
benchmark without any domain-specific fine-tuning. Medprompt is a fixed
composition of prompting techniques, not a search over prompts, and like most
prompt-optimisation methods it was evaluated on multiple-choice questions.

\begin{rwgap}
Prompt-optimisation methods are mostly evaluated on classification and
multiple-choice tasks with very large models. They have not been studied for
generating \emph{structured} patient-education content, such as video
storyboards, with a small, domain-fine-tuned model whose medical content must
not be changed by the optimised prompt.
\end{rwgap}

% ---------------------------------------------------------------------
\section{Automatic Generation of Educational Videos}
\label{sec:rw-video}
% ---------------------------------------------------------------------

Video is an effective format for patient education: a systematic review
found strong evidence that video-based interventions improve the knowledge
of patients with chronic illnesses~\cite{deshpande2023video}. Producing such
videos, however, requires medical experts, script writers and designers.
Recent work therefore investigates how LLMs can automate parts of this
process. TheoremExplainAgent~\cite{ku2025theoremexplainagent} uses LLM agents
to plan and write Manim animation code for explanation videos of
mathematical and scientific theorems; the best configuration produced
videos for 93.8\,\% of 240 theorems, although most videos still contained
minor visual layout problems. Code2Video~\cite{chen2026code2video} also
generates educational videos as executable code, using a planner, a coder and
a critic agent, and reports a 40\,\% improvement over direct code generation.
In healthcare, Haider et al.~\cite{haider2026ai} compared AI-generated avatar
videos with text handouts for postoperative education in a pilot randomised
trial. Knowledge scores did not differ significantly, but participants rated
the videos as clearer and more engaging.

These systems focus on the \emph{production} of a video from a given topic
and rely on the general knowledge of the LLM for its content. They do not
verify the medical statements against a trusted knowledge source, and they do
not adapt the content to the needs of a specific patient population.

\begin{rwgap}
LLM-based video generation has been studied mainly for STEM topics, while
medical studies use generic AI tools without knowledge grounding. A method is
missing that produces patient-oriented video storyboards whose medical
content is derived from guideline-based knowledge and whose structure can be
validated before production.
\end{rwgap}

% ---------------------------------------------------------------------
\section{Summary and Research Gap}
\label{sec:rw-summary}
% ---------------------------------------------------------------------

Table~\ref{tab:rw-comparison} compares representative approaches along the
dimensions that are relevant for this thesis. Each line of research covers
only part of the problem: medical LLMs and patient-education studies provide
patient-oriented language but no explicit knowledge; biomedical KG-RAG
systems ground answers in knowledge graphs but only at inference time;
neuro-symbolic methods learn from rules but have not been applied to
medicine; prompt optimisation and video generation have not been combined
with knowledge-grounded medical content.

\begin{table}[tbp]
  \centering
  \small
  \setlength{\tabcolsep}{4pt}
  \begin{tabular}{@{}lccccccc@{}}
    \toprule
    \textbf{Approach} & \textbf{Med.} & \textbf{KG} & \textbf{Rules}
      & \textbf{CoT-FT} & \textbf{P-Opt} & \textbf{Patient} & \textbf{Video} \\
    \midrule
    Med-PaLM~\cite{singhal2023large}                 & \rwyes & \rwno  & \rwno  & \rwno  & \rwno  & \rwpart & \rwno \\
    King et al.~\cite{king2025improving}             & \rwyes & \rwno  & \rwno  & \rwno  & \rwno  & \rwyes  & \rwno \\
    MedRAG~\cite{xiong2024benchmarking}              & \rwyes & \rwno  & \rwno  & \rwno  & \rwno  & \rwno   & \rwno \\
    KG-RAG~\cite{soman2024biomedical}                & \rwyes & \rwyes & \rwno  & \rwno  & \rwno  & \rwno   & \rwno \\
    DR.KNOWS~\cite{gao2025leveraging}                & \rwyes & \rwyes & \rwno  & \rwno  & \rwno  & \rwno   & \rwno \\
    MedGraphRAG~\cite{wu2025medical}                 & \rwyes & \rwyes & \rwno  & \rwno  & \rwno  & \rwno   & \rwno \\
    RoG~\cite{luo2024reasoning}                      & \rwno  & \rwyes & \rwpart& \rwpart& \rwno  & \rwno   & \rwno \\
    ChatRule~\cite{luo2025chatrule}                  & \rwno  & \rwyes & \rwyes & \rwno  & \rwno  & \rwno   & \rwno \\
    Chain-of-Knowledge~\cite{zhang2024chain}         & \rwno  & \rwyes & \rwyes & \rwyes & \rwno  & \rwno   & \rwno \\
    Medprompt~\cite{nori2023can}                     & \rwyes & \rwno  & \rwno  & \rwno  & \rwpart& \rwno   & \rwno \\
    GAAPO~\cite{secheresse2025gaapo}                 & \rwno  & \rwno  & \rwno  & \rwno  & \rwyes & \rwno   & \rwno \\
    TheoremExplainAgent~\cite{ku2025theoremexplainagent} & \rwno & \rwno & \rwno & \rwno & \rwno & \rwno   & \rwyes \\
    Code2Video~\cite{chen2026code2video}             & \rwno  & \rwno  & \rwno  & \rwno  & \rwno  & \rwno   & \rwyes \\
    \midrule
    \textbf{This thesis}                             & \rwyes & \rwyes & \rwyes & \rwyes & \rwyes & \rwyes  & \rwyes \\
    \bottomrule
  \end{tabular}
  \caption[Comparison of related approaches]{Comparison of related
  approaches. \textbf{Med.}: medical domain; \textbf{KG}: uses a knowledge
  graph; \textbf{Rules}: uses symbolic rules; \textbf{CoT-FT}: fine-tunes the
  LLM on (rule-based) reasoning chains; \textbf{P-Opt}: automatic prompt
  optimisation; \textbf{Patient}: patient-oriented output; \textbf{Video}:
  educational video or storyboard generation. \rwyes~=~yes,
  \rwpart~=~partially, \rwno~=~no. For this thesis, \textbf{Video} refers to
  storyboard scripts as input for video production.}
  \label{tab:rw-comparison}
\end{table}

Based on this analysis, the research gap addressed in this thesis can be
summarised as follows.

\begin{rwgap}[Research Gap of this Thesis]
There is no approach that (i)~represents clinical guidelines and patient
questionnaires about cardiovascular disease as knowledge graphs,
(ii)~turns symbolic rules mined from these graphs into chain-of-thought
training data for a medical LLM, (iii)~evaluates whether this improves
reasoning when knowledge is incomplete, and (iv)~uses the resulting model,
together with an automatically optimised prompt, to generate
patient-oriented storyboards for educational videos. This thesis addresses
this gap with the semantic-driven approach presented in
Chapter~4.
\end{rwgap}
